\documentclass[11pt,a4paper]{article}
\usepackage[T1]{fontenc}
\usepackage[utf8]{inputenc}
\usepackage{lmodern}
\usepackage{amsmath,amssymb,amsthm,mathtools}
\usepackage[margin=27mm,headheight=14pt]{geometry}
\usepackage{microtype}
\usepackage{enumitem}
\usepackage{booktabs}
\usepackage{xcolor}
\usepackage{fancyhdr}
\usepackage{xurl}
\usepackage{hyperref}
\usepackage{bookmark}
\usepackage[capitalise,nameinlink,noabbrev]{cleveref}
\definecolor{linkblue}{RGB}{22,68,105}
\hypersetup{colorlinks=true,linkcolor=linkblue,citecolor=linkblue,urlcolor=linkblue,
  pdftitle={A lower bound for the excess in the centered Poisson binomial comparison},
  pdfauthor={Yitzchak Shmalo},
  pdfsubject={Poisson binomial distribution, Kolmogorov distance, Hoeffding's interval inequality}}
\urlstyle{same}
\setlength{\parindent}{1.2em}
\setlength{\parskip}{0.2em}
\setlength{\emergencystretch}{1.5em}
\numberwithin{equation}{section}
\newtheorem{theorem}{Theorem}[section]
\newtheorem{lemma}[theorem]{Lemma}
\newtheorem{corollary}[theorem]{Corollary}
\newtheorem{proposition}[theorem]{Proposition}
\theoremstyle{remark}
\newtheorem{remark}[theorem]{Remark}
\newcommand{\Pp}{\mathbb P}
\newcommand{\E}{\mathbb E}
\newcommand{\Bin}{\operatorname{Bin}}
\pagestyle{fancy}
\fancyhf{}
\fancyhead[L]{\small Excess in the centered Poisson binomial comparison}
\fancyhead[R]{\small 19 November 2026}
\fancyfoot[C]{\thepage}
\renewcommand{\headrulewidth}{0.4pt}

\title{A lower bound for the excess in the centered Poisson binomial comparison}
\author{Yitzchak Shmalo\\[0.4em]
\small Einstein Institute of Mathematics, The Hebrew University of Jerusalem\\
\small Givat Ram, Jerusalem, Israel}
\date{19 November 2026}

\begin{document}
\maketitle
\thispagestyle{fancy}

\begin{abstract}
Let $X$ and $Y$ be sums of $n$ independent Bernoulli variables with parameters $p_i\ge q_i$, and let $g$ be the mean of the gaps $p_i-q_i$. Aryeh Kontorovich conjectured that for even $n$ the Kolmogorov distance between $X$ and $Y$ is at least $b_n(g)$, the distance between the binomial laws $\Bin(n,(1+g)/2)$ and $\Bin(n,(1-g)/2)$. Christopher D. Long proved this for every such pair. For centered pairs, those with $\sum_i(p_i+q_i)=n$, we bound the excess of the distance over $b_n(g)$ from below in two cases. When all gaps equal $g\in(0,1)$ and $n\ge4$ is even, both sums are built from one Bernoulli sum $W$ with mean $n/2$. Hoeffding's lower bound for the probability that $W$ lies in an interval centered at its mean then gives an excess of at least a positive multiple of $\sum_i(p_i+q_i-1)^2$. For $n=4$ and arbitrary gaps we give an exact formula for the excess of $\max(D_2,D_3)$, where $D_k=\Pp(X\ge k)-\Pp(Y\ge k)$.
\end{abstract}

\noindent\textit{Keywords:} Poisson binomial distribution, Kolmogorov distance, sums of independent Bernoulli variables, Hoeffding's inequalities for the number of successes.\\
\textit{2020 Mathematics Subject Classification:} 60E15, 60C05.

\section{Introduction}\label{sec:intro}

Let $X=X_1+\dots+X_n$ and $Y=Y_1+\dots+Y_n$, where the $X_i$ are independent Bernoulli$(p_i)$ variables and the $Y_i$ are independent Bernoulli$(q_i)$ variables, and suppose $p_i\ge q_i$ for every $i$. The laws of $X$ and $Y$ are Poisson binomial laws, and only these laws matter, so no coupling between the two sums is assumed. Write
\[
d_i=p_i-q_i,\qquad s_i=p_i+q_i-1,\qquad \Delta=\sum_i d_i,\qquad g=\frac{\Delta}{n},\qquad S_2=\sum_i s_i^2,
\]
so that the $d_i$ are the gaps and $g$ is the mean gap. For every $n$ we also write
\begin{equation}\label{eq:Qdef}
Q=\sum_i\Bigl(d_i-\frac\Delta2\Bigr)s_i^2 .
\end{equation} Call $D_k=\Pp(X\ge k)-\Pp(Y\ge k)$ the tail difference at $k$. A coordinatewise coupling shows $D_k\ge0$, so the Kolmogorov distance is $d_K(X,Y)=\max_k D_k$. For even $n$, writing $\Bin(n,p)$ for a binomial variable with $n$ trials and success probability $p$, put
\[
b_n(g)=\Pp\{\Bin(n,\tfrac{1+g}2)\ge\tfrac n2\}-\Pp\{\Bin(n,\tfrac{1-g}2)\ge\tfrac n2\}.
\]
In a question on MathOverflow~\cite{Kontorovich}, Aryeh Kontorovich conjectured that
\begin{equation}\label{eq:conj}
d_K(X,Y)\ge b_n(g)\qquad\text{for every even } n.
\end{equation}
The right side is the Kolmogorov distance of the symmetric binomial pair $\Bin(n,\frac{1+g}2)$, $\Bin(n,\frac{1-g}2)$, which has the same mean gap; it is attained at $k=n/2$. Since $p_i=(1+s_i+d_i)/2$ and $q_i=(1+s_i-d_i)/2$, both lie in $[0,1]$ precisely when $|s_i|+d_i\le1$. Write $s=(s_1,\dots,s_n)$ and $d=(d_1,\dots,d_n)$; we call $(s,d)$, and the pair $X,Y$, feasible when $|s_i|+d_i\le1$ for every $i$. Every pair $X,Y$ is feasible, and every feasible $(s,d)$ comes from a pair.

Christopher D. Long~\cite[Theorem~1.1]{Long} proved~\eqref{eq:conj} for every such pair, together
with its analogue for odd $n$; following him we call~\eqref{eq:conj} the Poisson binomial comparison, and for a quantity $T\le d_K(X,Y)$ we call $T-b_n(g)$ the excess of $T$. He also showed that for even $n\ge4$ and $0<g<1$ equality holds only for
the symmetric binomial pair. His theorem determines the pairs with $d_K(X,Y)=b_n(g)$ and gives no
lower bound for $d_K(X,Y)-b_n(g)$ at the other pairs, that is, no quantitative form of the equality statement. The proof is an induction on $n$: at a fixed mean gap it takes a pair minimizing $\max_kD_k$, which exists by compactness, and identifies it by a stationary-point analysis that includes the boundary faces. The argument locates the minimizer and does not estimate how fast $\max_kD_k$ grows away from it. Call the pair centered when $\sum_is_i=0$,
that is, $\E X+\E Y=n$, and write $n=2m$. For a sum $S$ of $n$ independent Bernoulli variables and an
integer $c$, Hoeffding's tail comparison~\cite[Theorem~4]{Hoeffding}, which holds in the form with non-strict
inequalities in the conclusion~\cite[Theorem~2.1(1)]{TangTang}, states that $\Pp(S\le c)$ is at most the
corresponding probability for the binomial law with the same mean when $c\le\E S-1$, and at least
that probability when $c\ge\E S$. For a centered pair, $\E X=m+\frac{ng}2$ and
$\E Y=m-\frac{ng}2$, and the two binomial laws are those of the symmetric binomial pair. If
$g\ge2/n$, then $\E Y\le m-1\le\E X-1$, and the comparison at $c=m-1$, applied to $X$ and to $Y$ as
in Sawin's answer~\cite{Sawin}, gives $D_m\ge b_{2m}(g)$. This argument uses the single threshold $k=m$, requires
$g\ge2/n$, and gives no term beyond $b_{2m}(g)$.

Suppose now that $n=2m\ge4$, the pair is centered and every gap equals $g$, with $0<g<1$. Let $R_1,\dots,R_n$ be independent Bernoulli
variables with parameters $r_i=\frac12+\frac{s_i}{2(1-g)}$, let $Z_1,\dots,Z_n$ be independent
Bernoulli$(g)$ variables independent of the $R_i$, and put $W=R_1+\dots+R_n$. Then
$\sum_i\max(R_i,Z_i)$ has the law of $X$, $\sum_iR_i(1-Z_i)$ has the law of $Y$, and $\E W=m$.
Given $W=w$, the two sums have the laws of $w+\Bin(2m-w,g)$ and $\Bin(w,1-g)$, and the value of
$(D_m+D_{m+1})/2$ for that pair is symmetric about $w=m$ and decreases as $|w-m|$ grows.
Hoeffding's interval inequality~\cite[Theorem~5]{Hoeffding} gives
\[
\Pp(|W-m|\le j)\ \ge\ \Pp\bigl\{|\Bin(2m,\tfrac12)-m|\le j\bigr\},\qquad 0\le j\le m-1 .
\]
Hence $(D_m+D_{m+1})/2-b_{2m}(g)$ is a combination, with positive coefficients, of the differences
of the two sides of this inequality, and these differences, weighted by $2j+1$, sum to
$\operatorname{Var}\Bin(2m,\tfrac12)-\operatorname{Var}W=S_2/(4(1-g)^2)$. Comparing each coefficient with $2j+1$ gives \cref{thm:equal} below, and since
$d_K(X,Y)\ge(D_m+D_{m+1})/2$ its bound holds for $d_K(X,Y)$ as well. For $m=2$ the bound of \cref{thm:equal} is an equality, and for $m\ge3$ and small $g$ there is a pair at which the excess of $(D_m+D_{m+1})/2$ is much larger than the last term of~\eqref{eq:equal} (\cref{rem:size}). For $n=4$ and arbitrary gaps,
\cref{thm:four} gives $\max(D_2,D_3)-b_4(g)$ exactly.

\begin{theorem}[equal gaps]\label{thm:equal}
Let $n=2m\ge4$, let the pair be feasible, suppose $d_i=g$ for every $i$, and suppose $\sum_i s_i=0$. Then
\begin{equation}\label{eq:equal}
\frac{D_m+D_{m+1}}2\ge b_{2m}(g)+\frac{m+1}{24}\,g^{m-1}(1-g)^{m-2}S_2 .
\end{equation}
In particular $d_K(X,Y)\ge b_{2m}(g)$, and for $g>0$ equality $d_K(X,Y)=b_{2m}(g)$ holds only when every $s_i$ vanishes.
\end{theorem}

The last two statements are cases of \cite[Theorem~1.1]{Long}; the new part of \cref{thm:equal} is the term in $S_2$. The excess of $(D_m+D_{m+1})/2$ over $b_{2m}(g)$ has an exact expression as a sum of nonnegative terms, given in~\eqref{eq:exact}, and~\eqref{eq:cstar} gives a coefficient of $S_2$ at least as large as the one in~\eqref{eq:equal}. For $n=2$ the average $(D_1+D_2)/2$ equals $b_2(g)=g$ for every pair (\cref{rem:n2}), so no positive multiple of $S_2$ can be added there.

\begin{theorem}[four summands]\label{thm:four}
Let $n=4$, let the pair be feasible with $\sum_i s_i=0$, and let the gaps $d_i\ge0$ be arbitrary. Let $Q$ be as in~\eqref{eq:Qdef} and put
\[
L=\sum_i s_id_i+3\sum_i\Bigl(d_i\prod_{j\ne i}s_j+s_i\prod_{j\ne i}d_j\Bigr).
\]
Then $Q\le|L|$, and $|L|\ge\frac83Q$ when $Q>0$; the constant $\frac83$ cannot be increased. Moreover
\begin{equation}\label{eq:four}
\max(D_2,D_3)-b_4(g)=\frac1{384}\sum_{i<j}(d_i-d_j)^2\bigl[5\Delta-4(d_i+d_j)\bigr]+\frac{|L|-Q}8\ \ge0 ,
\end{equation}
and hence $d_K(X,Y)\ge\max(D_2,D_3)\ge b_4(g)$. Since $|L|-Q\ge|Q|$, the excess of $\max(D_2,D_3)$ is at least the first term of~\eqref{eq:four} plus $|Q|/8$.
For $\Delta>0$, equality $d_K(X,Y)=b_4(g)$ holds only for the symmetric binomial pair.
\end{theorem}

The inequality $d_K(X,Y)\ge b_4(g)$ and its equality statement are cases of \cite[Theorem~1.1]{Long}. When every gap equals $g$, the first term of~\eqref{eq:four} vanishes and $|Q|/8=gS_2/8$, the bound of \cref{thm:equal} for $m=2$.

Two answers on MathOverflow contain partial results. Will Sawin~\cite{Sawin} uses an exchange argument. He shows that if an integer $k$ satisfies $\E Y+1<k<\E X$, then $\Pp(X\ge k)$ is at least, and $\Pp(Y\ge k)$ at most, the corresponding tail of the binomial law with the same mean, a case of Hoeffding's tail comparison stated above. For pairs that are not centered the two binomial laws have success probabilities $\E X/n$ and $\E Y/n$, so that the comparison of the difference of their tails with $b_n(g)$ is a further step. Leon Jay Gleser~\cite{Gleser} sharpens Hoeffding's tail comparison to Schur-convexity statements for the distribution function of one sum; they give no bound for the difference of the binomial tails with success probabilities $\E X/n$ and $\E Y/n$ in terms of $b_n(g)$.

Thomas Steinke~\cite{Steinke} considers the total variation distance between the laws of the vectors $(X_1,\dots,X_n)$ and $(Y_1,\dots,Y_n)$, which are product measures. This distance is at least $d_K(X,Y)$, because $\{X\ge k\}$ is an event of the product space, and for the symmetric binomial pair the likelihood ratio is an increasing function of the sum, so the two distances coincide and equal $b_n(g)$. Hence~\eqref{eq:conj} implies the inequality in which $d_K(X,Y)$ is replaced by the total variation distance. This inequality, with its equality cases, is part of \cite[Theorem~1.1]{Long}. Earlier, Steinke proved it up to a factor of two, by mapping each pair of Bernoulli laws to a pair symmetric about $\frac12$ with a randomized function.

Wenpin Tang and Fengmin Tang~\cite{TangTang} survey Hoeffding's and Gleser's inequalities and other properties of Poisson binomial laws. Kontorovich~\cite{KontorovichTens} bounds the total variation distance between two product measures from below in terms of the distances between their factors. In~\cite{KontorovichHom,KontorovichTV} he proves that the total variation distance between two products of Bernoulli laws, and more generally of arbitrary laws, is at least a universal constant times the distance between the products in which every factor is replaced by the average of its family. For Bernoulli products the best such constant is at most $8/9$~\cite[remark after Theorem~1.5]{KontorovichTV}. For centered pairs these averaged products are the symmetric binomial pair, so~\eqref{eq:conj} asserts for such pairs the corresponding inequality for the Kolmogorov distance with constant~$1$. Since $d_K(X,Y)$ is at most the total variation distance, these lower bounds do not bound $d_K(X,Y)$ from below.

The lower bound in~\eqref{eq:equal} is for the average of $D_m$ and $D_{m+1}$, which is at most $d_K(X,Y)$; \cite[Theorem~1.1]{Long} is stated for the maximum. For $m\ge3$ the largest $\gamma$ such that $(D_m+D_{m+1})/2\ge b_{2m}(g)+\gamma S_2$ for all centered pairs with equal gaps is not known, and no positive lower bound for $d_K(X,Y)-b_n(g)$ in terms of the $s_i$ and $d_i$ is known for centered pairs with unequal gaps and $n\ge6$. For $n=6$ and centered pairs, $\max(D_3,D_4)-b_6(g)$ is again the sum of two terms: the first depends on the gaps alone and is nonnegative, and it is not known whether the second is nonnegative for every centered pair. If it is, then $\max(D_3,D_4)\ge b_6(g)$ for every centered pair with $n=6$ (\cref{sec:six}).

\section{Equal gaps}\label{sec:equal}

We use Hoeffding's interval inequality~\cite[Theorem~5]{Hoeffding}: if $S$ is a sum of $n$ independent Bernoulli variables with $\E S=n\bar r$, so that $\bar r$ is the mean of their parameters, and $\ell_1,\ell_2$ are integers with $0\le \ell_1\le n\bar r\le \ell_2\le n$, then
\begin{equation}\label{eq:hoeffding}
\Pp(\ell_1\le S\le \ell_2)\ge\Pp\{\ell_1\le\Bin(n,\bar r)\le \ell_2\}.
\end{equation}
We apply it with $n\bar r=m$ an integer and the intervals $[m-j,m+j]$, $0\le j\le m-1$.

\begin{proof}[Proof of \cref{thm:equal}]
If $g=0$ then $p_i=q_i$ for every $i$ and there is nothing to prove; if $g=1$ then feasibility forces $p_i=1$ and $q_i=0$, so $S_2=0$ and both sides equal $1$. Let $0<g<1$, put $a=1-g$, and let $r_i$, $R_i$, $Z_i$ and $W$ be as in \cref{sec:intro}, so that $r_i=(1+s_i/a)/2$. Feasibility $|s_i|\le a$ gives $r_i\in[0,1]$, centering gives $\sum_i r_i=m$, and $q_i=ar_i$, $p_i=g+ar_i$.

Put
\[
\widetilde X_i=\max(R_i,Z_i),\qquad \widetilde Y_i=R_i(1-Z_i).
\]
Then $\widetilde X_i$ is Bernoulli$(p_i)$ and $\widetilde Y_i$ is Bernoulli$(q_i)$, with independent coordinates in each family, so $\sum_i\widetilde X_i$ has the law of $X$ and $\sum_i\widetilde Y_i$ the law of $Y$. Given $W=w$, the first sum is $w$ plus the number of indices with $R_i=0$ and $Z_i=1$, and the second is the number of indices with $R_i=1$ and $Z_i=0$. Since $\Pp(Z_i=1)=g$ is the same in every coordinate, with $\overset d=$ denoting equality in law,
\begin{equation}\label{eq:conditional}
\Bigl(\sum_i\widetilde X_i\Bigm| W=w\Bigr)\overset d=w+\Bin(2m-w,g),\qquad
\Bigl(\sum_i\widetilde Y_i\Bigm| W=w\Bigr)\overset d=\Bin(w,a).
\end{equation}

Let $h(k)=-1,0,1$ according as $k<m$, $k=m$, $k>m$, so that $h(k)=\mathbf 1_{\{k\ge m\}}+\mathbf 1_{\{k\ge m+1\}}-1$ with $\mathbf 1$ denoting an indicator. Define
\[
\phi(w)=\tfrac12\Bigl(\E h\bigl(w+\Bin(2m-w,g)\bigr)-\E h\bigl(\Bin(w,a)\bigr)\Bigr),\qquad 0\le w\le 2m .
\]
By~\eqref{eq:conditional},
\begin{equation}\label{eq:average}
\frac{D_m+D_{m+1}}2=\E\phi(W).
\end{equation}
The function $\phi$ is symmetric about $m$. Indeed $2m-(w+\Bin(2m-w,g))$ has the law of $\Bin(2m-w,a)$, $2m-\Bin(w,a)$ has the law of $(2m-w)+\Bin(w,g)$, and $h(2m-k)=-h(k)$; hence $\phi(w)=\phi(2m-w)$.

We compute $\phi$ on $[0,m]$. If $w=m-j$ with $1\le j\le m$, then $\Bin(m-j,a)<m$ always, and with $B$ a $\Bin(m+j,g)$ variable,
\[
\phi(m-j)=\Pp(B>j)+\tfrac12\Pp(B=j).
\]
At the center, $\phi(m)=1-a^m$. Let $A_j=\phi(m-j)-\phi(m-j-1)$, $0\le j\le m-1$, be the jumps of $\phi$. For $j\ge1$, let $B'=B+\xi$ with $\xi$ an independent Bernoulli$(g)$ variable, so that $\phi(m-j-1)=\Pp(B'>j+1)+\frac12\Pp(B'=j+1)$. Expanding $B'$ in terms of $B$ gives $A_j=\frac a2\bigl(\Pp(B=j)+\Pp(B=j+1)\bigr)$, and direct substitution at $j=0$ gives $A_0$:
\begin{equation}\label{eq:jumps}
A_0=\frac{m-1}2\,g\,a^m,\qquad
A_j=\frac12\binom{m+j}{j}g^ja^m\Bigl(a+\frac{m}{j+1}\,g\Bigr)\quad(1\le j\le m-1).
\end{equation}
All $A_j$ are positive when $m\ge2$. Telescoping and symmetry give, for $0\le w\le 2m$,
\begin{equation}\label{eq:expansion}
\phi(w)=\phi(0)+\sum_{j=0}^{m-1}A_j\,\mathbf 1_{\{|w-m|\le j\}} .
\end{equation}

Let $B_0$ have the law $\Bin(2m,\frac12)$, and put $\delta_j=\Pp(|W-m|\le j)-\Pp(|B_0-m|\le j)$ for $0\le j\le m-1$. By~\eqref{eq:hoeffding} with $\ell_1=m-j$ and $\ell_2=m+j$, every $\delta_j$ is nonnegative. When every $r_i$ equals $\frac12$, $\sum_i\widetilde X_i$ and $\sum_i\widetilde Y_i$ form the symmetric binomial pair, $W$ has the law of $B_0$, and~\eqref{eq:average} equals $(D_m+D_{m+1})/2$ for that pair. By the symmetry $\Bin(2m,\frac{1-g}2)\overset d=2m-\Bin(2m,\frac{1+g}2)$ this is $b_{2m}(g)$. Taking expectations in~\eqref{eq:expansion} under $W$ and under $B_0$ gives
\begin{equation}\label{eq:exact}
\frac{D_m+D_{m+1}}2-b_{2m}(g)=\sum_{j=0}^{m-1}A_j\delta_j\ \ge0 .
\end{equation}

To bound the right side of~\eqref{eq:exact}, note that an integer $\nu\in\{0,\dots,m\}$ satisfies $\nu^2=\sum_{j=0}^{m-1}(2j+1)\mathbf 1_{\{\nu>j\}}$. Applying this to $|W-m|$ and $|B_0-m|$, which take values in $\{0,\dots,m\}$, and using $\E W=\E B_0=m$,
\begin{equation}\label{eq:variance}
\sum_{j=0}^{m-1}(2j+1)\delta_j=\operatorname{Var}B_0-\operatorname{Var}W=\sum_i\Bigl(r_i-\frac12\Bigr)^2=\frac{S_2}{4a^2}.
\end{equation}
Hence~\eqref{eq:exact} is at least $c^*_{2m}(g)\,S_2$ with
\begin{equation}\label{eq:cstar}
c^*_{2m}(g)=\frac1{4a^2}\min_{0\le j\le m-1}\frac{A_j}{2j+1}.
\end{equation}
For $1\le j\le m-1$ we have $a+mg/(j+1)\ge a+g=1$ and $g^j\ge g^{m-1}$. The ratio $\binom{m+j}{j}/(2j+1)$ is nondecreasing in $j\ge1$, because
\[
(m+j+1)(2j+1)-(j+1)(2j+3)=(2m-2)j+m-2\ge0,
\]
so it is at least its value $(m+1)/3$ at $j=1$. For $j=0$ we use $m-1\ge(m+1)/3$ and $g\ge g^{m-1}$. Therefore $A_j/(2j+1)\ge\frac{m+1}6g^{m-1}a^m$ for every $j$, and~\eqref{eq:cstar} is at least $\frac{m+1}{24}g^{m-1}a^{m-2}$. This proves~\eqref{eq:equal}.

Since $d_K(X,Y)\ge(D_m+D_{m+1})/2$, we get $d_K(X,Y)\ge b_{2m}(g)$, with equality for $0<g<1$ only if $S_2=0$. Conversely, for the symmetric binomial pair the ratio $\Pp(X=k)/\Pp(Y=k)=\bigl((1+g)/(1-g)\bigr)^{2k-n}$ is below $1$ for $k<m$ and above $1$ for $k>m$, so $D_k$ is largest at $k=m$ and $d_K(X,Y)=b_{2m}(g)$.
\end{proof}

\begin{remark}\label{rem:n2}
For $n=2$ and every pair, $D_1+D_2=\E X-\E Y=2g$, while both tail differences of the symmetric binomial pair equal $\bigl((1+g)^2-(1-g)^2\bigr)/4=g$, so $b_2(g)=g$. Hence $d_K(X,Y)\ge(D_1+D_2)/2=b_2(g)$, which is the case $n=2$ of \cite[Theorem~1.1]{Long}, without centering and with arbitrary gaps, and the average leaves no room for a positive multiple of $S_2$. In the construction above this is the case $m=1$, where the jump $A_0$ vanishes and $\phi$ is constant.
\end{remark}

\begin{remark}\label{rem:size}
For $m=2$ one has $A_0=A_1/3$, so~\eqref{eq:exact} and~\eqref{eq:variance} give $(D_2+D_3)/2-b_4(g)=\frac g8S_2$, and~\eqref{eq:equal} is an equality. For every $m\ge2$,~\eqref{eq:exact} and~\eqref{eq:variance} also give
\[
\frac{D_m+D_{m+1}}2-b_{2m}(g)\le\frac1{4a^2}\max_{0\le j\le m-1}\frac{A_j}{2j+1}\,S_2 ,
\]
so for fixed $m$ and $0<g<1$ this excess lies between two positive multiples of $S_2$.
Let $m\ge3$, and let $s_i=1-g$ for $m$ indices and $s_i=-(1-g)$ for the others. Then $W=m$ almost surely, and~\eqref{eq:exact} gives $\frac{D_m+D_{m+1}}2-b_{2m}(g)\ge A_0\delta_0=\frac{m-1}2g(1-g)^m\bigl(1-\binom{2m}m4^{-m}\bigr)$, whereas the last term of~\eqref{eq:equal} equals $\frac{m(m+1)}{12}g^{m-1}(1-g)^m$. The quotient of the two is $\frac{6(m-1)}{m(m+1)}\bigl(1-\binom{2m}m4^{-m}\bigr)g^{-(m-2)}$, so for $m\ge3$ and small $g$ the bound~\eqref{eq:equal} is far from the excess of $(D_m+D_{m+1})/2$ at this pair.
\end{remark}

\begin{remark}\label{rem:b2m}
The value $b_{2m}(g)$ is explicit: differentiating the two binomial tails gives
\[
b_{2m}'(g)=\frac{2m}{2^{2m-1}}\binom{2m-1}{m-1}(1-g^2)^{m-1},\qquad b_{2m}(0)=0,
\]
so that, for example, $b_4(g)=\frac32g-\frac12g^3$.
\end{remark}

\section{Four summands with arbitrary gaps}\label{sec:four}

The proof of \cref{thm:four} reduces the inequality $Q\le|L|$ to the following statement about three numbers.

\begin{lemma}\label{lem:three}
Let $-\frac13\le y_i\le1$ for $i=1,2,3$ with $y_1+y_2+y_3=1$, and let $v_i\ge0$ with $v_1+v_2+v_3=1$. Put
\[
U=\sum_i v_i(1-y_i),\qquad F=\sum_i y_i\prod_{j\ne i}v_j-v_1v_2v_3 .
\]
Then $U\ge0$ and $3F\le U$.
\end{lemma}

\begin{proof}
Each term of $U$ is nonnegative. Both $U$ and $F$ are unchanged when the indices of $y$ and $v$ are permuted together. For fixed $v$ the quantity $U-3F$ is affine in $y$, and the polygon $\{y\in[-\frac13,1]^3:\sum_iy_i=1\}$ has as vertices the six permutations of $(-\frac13,\frac13,1)$. So it suffices to take $y=(-\frac13,\frac13,1)$ and an arbitrary $v$. Multiplying $U$ by $(v_1+v_2+v_3)^2=1$ and the quadratic part of $F$ by $v_1+v_2+v_3=1$ makes both cubic, and a direct expansion gives
\[
3(U-3F)=4v_1^3+v_2\Bigl(v_1-\frac{v_2}2\Bigr)^2+\frac74v_2^3+5v_1^2v_3+7v_2^2v_3+v_1v_3^2+5v_2v_3^2+12v_1v_2v_3,
\]
in which every term is nonnegative.
\end{proof}

\begin{proof}[Proof of \cref{thm:four}]
For $z\in\{-1,1\}^4$ let $M(z)$ be $1$, $0$ or $-1$ according as the number of positive coordinates is larger than, equal to or smaller than two, and let $e_k$ be the $k$-th elementary symmetric polynomial of its arguments. Checking the five possible numbers of positive signs shows $M=(3e_1-e_3)/8$. If $z_i=2X_i-1$, then $M(z)=\mathbf 1_{\{X\ge2\}}+\mathbf 1_{\{X\ge3\}}-1$, and similarly for $Y$; hence $D_2+D_3=\E M(2X_\bullet-1)-\E M(2Y_\bullet-1)$, where $2X_\bullet-1$ denotes the vector $(2X_i-1)_i$. The coordinates $2X_i-1$ are independent with means $s_i+d_i$, and the coordinates $2Y_i-1$ are independent with means $s_i-d_i$. Since $M$ is multilinear, each expectation is $(3e_1-e_3)/8$ evaluated at the vector of means. With $\sum_is_i=0$ one has $e_2(s)=-S_2/2$, and the part of $e_3(s+d)$ that is odd in $d$ equals $e_3(d)+\sum_id_i\bigl(s_i^2-S_2/2\bigr)=e_3(d)+Q$. Therefore
\begin{equation}\label{eq:avg4}
\frac{D_2+D_3}2=\frac{3\Delta-e_3(d)-Q}8 .
\end{equation}
Since $\Pp(X=2)=e_2(p)-3e_3(p)+6e_4(p)$ and likewise for $Y$, the substitution $p_i=(1+s_i+d_i)/2$, $q_i=(1+s_i-d_i)/2$ gives
\begin{equation}\label{eq:diff4}
D_2-D_3=\Pp(X=2)-\Pp(Y=2)=\frac L4 .
\end{equation}
Hence
\begin{equation}\label{eq:max4}
\max(D_2,D_3)=\frac{3\Delta-e_3(d)-Q+|L|}8 .
\end{equation}

We show $Q\le|L|$, and $|L|\ge\frac83Q$ when $Q>0$. There is nothing to prove if $Q\le0$, so let $Q>0$. Since $Q=\sum_id_i(s_i^2-S_2/2)$ and every $d_i\ge0$, some $s_i^2$ exceeds $S_2/2$, and at most one can. Relabel it as $s_1$. Replacing every $s_i$ by $-s_i$, which corresponds to replacing $(p_i,q_i)$ by $(1-q_i,1-p_i)$, keeps the pair feasible and centered, leaves $Q$ unchanged and changes the sign of $L$; so we may assume $s_1>0$ and write
\[
s=(x,-xy_1,-xy_2,-xy_3),\qquad d=(t,u_1,u_2,u_3),\qquad \sigma=u_1+u_2+u_3,
\]
with $x>0$ and $y_1+y_2+y_3=1$. Put $\rho=y_1^2+y_2^2+y_3^2$. The choice of $s_1$ means $x^2>S_2/2=x^2(1+\rho)/2$, that is, $\rho<1$. Since $y_i^2+(1-y_i)^2/2\le\rho<1$, each $y_i$ lies in $(-\frac13,1)$, and at most one $y_i$ is negative. Put $\kappa=(1-\rho)/2$ and $\theta_i=(1+\rho)/2-y_i^2$. Then $0<\kappa\le\frac13$ because $\rho\ge\frac13$, and $\theta_i-\kappa=\rho-y_i^2\ge0$. Expanding $Q$ in these coordinates and using $u_i\ge0$,
\begin{equation}\label{eq:Q}
Q=x^2\Bigl(\kappa t-\sum_i\theta_iu_i\Bigr)\le\kappa x^2(t-\sigma).
\end{equation}
Hence $t>\sigma\ge0$, and since $\kappa\le\frac13$, $Q\le x^2(t-\sigma)/3$. Feasibility at the first coordinate gives $x+t\le1$, so $0<x<1$ and $\sigma<1-x$.

As a function of $t$ with the other variables fixed, $L/x$ is affine; expanding $L$ gives
\begin{equation}\label{eq:affine}
\begin{aligned}
\frac Lx&=tC+G,\qquad C=1-3x^2y_1y_2y_3-3\sum_iy_i\prod_{j\ne i}u_j,\\
G&=\sum_iu_i\Bigl(3x^2\prod_{j\ne i}y_j-y_i\Bigr)+3u_1u_2u_3 .
\end{aligned}
\end{equation}
Because at most one $y_i$ is negative, $y_1y_2y_3\le\frac1{27}$, either by its sign or by the inequality of arithmetic and geometric means. Because each $y_i\le1$, $\sum_iy_i\prod_{j\ne i}u_j\le\sum_{i<j}u_iu_j\le\sigma^2/3$. With $\sigma<1-x$,
\begin{equation}\label{eq:slope}
C\ge1-\frac{x^2}9-\sigma^2\ge1-\frac{x^2}9-(1-x)^2=2x-\frac{10}9x^2\ge\frac89x .
\end{equation}

Next we show that $\sigma C+G$, the value of $L/x$ at $t=\sigma$, is nonnegative; the point $t=\sigma$ need not be feasible. If $\sigma=0$ then $G=0$. If $\sigma>0$ put $v_i=u_i/\sigma$, so that $v_1+v_2+v_3=1$, let $U$ and $F$ be as in \cref{lem:three}, and put $V=\sum_iv_i(1-y_i)\prod_{j\ne i}y_j$. Substituting the definitions of $C$ and $G$ and collecting terms gives
\begin{equation}\label{eq:regroup}
\begin{aligned}
\sigma C+G&=\sum_iu_i(1-y_i)+3x^2\sum_iu_i(1-y_i)\prod_{j\ne i}y_j-3\Bigl(\sigma\sum_iy_i\prod_{j\ne i}u_j-u_1u_2u_3\Bigr)\\
&=\sigma\bigl(U+3x^2V-3\sigma^2F\bigr).
\end{aligned}
\end{equation}
Each product $y_jy_k$ of two distinct $y$'s is at least $-\frac13$, since at most one factor is negative, that factor exceeds $-\frac13$, and the other is below $1$. As $v_i(1-y_i)\ge0$, this gives $V\ge-U/3$, and \cref{lem:three} gives $3F\le U$. Therefore
\begin{equation}\label{eq:atk}
U+3x^2V-3\sigma^2F\ge U(1-x^2-\sigma^2)\ge U\bigl(1-x^2-(1-x)^2\bigr)=2Ux(1-x)\ge0 .
\end{equation}
Combining~\eqref{eq:affine},~\eqref{eq:slope},~\eqref{eq:atk} and~\eqref{eq:Q},
\[
L=x\bigl(\sigma C+G+(t-\sigma)C\bigr)\ge\frac89x^2(t-\sigma)\ge\frac83Q .
\]
Since replacing $s$ by $-s$ changes only the sign of $L$, $|L|\ge\frac83Q$ whenever $Q>0$.

For sharpness take $d=(t,0,0,0)$ and $s=(x,-\frac x3,-\frac x3,-\frac x3)$ with $x=1-t$ and $0<t<1$. This pair is feasible and centered, $Q=tx^2/3>0$ and $L=xt-tx^3/9$, so $|L|/Q=3/x-x/3$, which tends to $\frac83$ as $t\to0$.

Finally, the identity
\begin{equation}\label{eq:gap}
\frac{\Delta^3}{16}-e_3(d)=\frac1{48}\sum_{i<j}(d_i-d_j)^2\bigl[5\Delta-4(d_i+d_j)\bigr],
\end{equation}
together with~\eqref{eq:max4} and $b_4(\Delta/4)=\frac38\Delta-\frac1{128}\Delta^3$ (\cref{rem:b2m}), gives~\eqref{eq:four}. Each bracket in~\eqref{eq:gap} equals $\Delta+4(\Delta-d_i-d_j)\ge\Delta$. So if $\Delta>0$ and $d_K(X,Y)=b_4(g)$, then $\max(D_2,D_3)=b_4(g)$, both terms of~\eqref{eq:four} vanish, and all gaps are equal. Then $Q=-gS_2$, so $|L|-Q\ge gS_2$ forces $S_2=0$. The symmetric binomial pair attains equality, by the last paragraph of the proof of \cref{thm:equal}.
\end{proof}

\begin{corollary}\label{cor:single}
If $n=4$, $\sum_is_i=0$ and only one gap is nonzero, say $d=(t,0,0,0)$, then $d_K(X,Y)\ge\frac38t=b_4(t/4)+\frac{t^3}{128}$.
\end{corollary}

\begin{proof}
In~\eqref{eq:four} the three pairs $(1,j)$ contribute $t^2(5t-4t)=t^3$ each and the other pairs contribute nothing, so the first term is $3t^3/384=t^3/128$.
\end{proof}

\section{Six summands}\label{sec:six}
For $n=6$ the function on $\{-1,1\}^6$ equal to the sign of the number of positive coordinates minus three is $M_6=(5e_1-e_3+e_5)/16$, and $D_3-D_4=\Pp(X=3)-\Pp(Y=3)$. With these in place of $M$ and $D_2-D_3$, for centered pairs the argument that gives~\eqref{eq:avg4} and~\eqref{eq:diff4} gives
\[
D_3-D_4=\frac{L_6}8,\qquad \frac{D_3+D_4}2=\frac{5\Delta-e_3(d)+e_5(d)-Q+H}{16},
\]
where $Q$ is given by~\eqref{eq:Qdef} with six summands and, writing $s^{(I)}$ for the vector $s$ with the coordinates indexed by $I$ removed,
\[
H=\sum_i d_i\,e_4\bigl(s^{(\{i\})}\bigr)+\sum_{|I|=3}\Bigl(\prod_{i\in I}d_i\Bigr)e_2\bigl(s^{(I)}\bigr),
\]
\[
L_6=\sum_i s_id_i+O_4-5O_6,\qquad O_j=\frac{e_j(s+d)-e_j(s-d)}2 .
\]
For the symmetric binomial pair $Q=H=L_6=0$ and $\Delta=6g$, so $b_6(g)=(30g-20g^3+6g^5)/16$.

\begin{proposition}\label{prop:six}
For every centered pair with $n=6$,
\begin{equation}\label{eq:six}
\max(D_3,D_4)-b_6(g)=\frac{20g^3-6g^5-e_3(d)+e_5(d)}{16}+\frac{|L_6|-(Q-H)}{16},
\end{equation}
and the first term on the right is nonnegative.
\end{proposition}

\begin{proof}
Since $\max(D_3,D_4)=(D_3+D_4)/2+|D_3-D_4|/2$ and $5\Delta=30g$, the two displays above and the value of $b_6(g)$ give~\eqref{eq:six}. Since $d_i\le1-|s_i|$, $d\in[0,1]^6$. For $\eta\in[0,1]^4$ we have $e_3(\eta)\le e_1(\eta)$, since the product of the three coordinates other than $\eta_k$ is at most $\eta_{k+1}$ (indices mod~$4$). If $d=(d_i,d_j,\eta)$ after a permutation, then
\[
e_3(d)-e_5(d)=d_id_j\bigl(e_1(\eta)-e_3(\eta)\bigr)+(d_i+d_j)\bigl(e_2(\eta)-e_4(\eta)\bigr)+e_3(\eta),
\]
so replacing $d_i$ and $d_j$ by their average does not decrease $e_3(d)-e_5(d)$. Averaging the largest and the smallest coordinate repeatedly converges to $(g,\dots,g)$, so by continuity $e_3(d)-e_5(d)\le20g^3-6g^5$.
\end{proof}

By \cref{prop:six}, the inequality
\begin{equation}\label{eq:open6}
Q-H\le|L_6|
\end{equation}
implies $\max(D_3,D_4)\ge b_6(g)$ for every centered pair with $n=6$. Neither~\eqref{eq:open6} nor the inequality $\max(D_3,D_4)\ge b_6(g)$ has been proved for every centered pair with $n=6$; \cite[Theorem~1.1]{Long} gives $d_K(X,Y)\ge b_6(g)$, where the maximum is over all $k$. Three examples bear on~\eqref{eq:open6}.

On the family $s=(\beta,-\beta,0,0,0,0)$, with $\tau=(d_3,\dots,d_6)$,
\[
Q-H=\beta^2\bigl(e_3(\tau)-e_1(\tau)\bigr)\le0,\qquad L_6=\beta(d_1-d_2)\bigl(1-e_2(\tau)+5e_4(\tau)\bigr),
\]
and both vanish when $\tau=(1,1,1,1)$, or when $\tau=0$ and $d_1=d_2$. So~\eqref{eq:open6} holds with equality, $Q-H=|L_6|=0$, at feasible pairs with $s\ne0$.

In \cref{sec:four}, $Q>0$ forces some $s_i^2$ to exceed $S_2/2$. For $n=6$, $Q-H$ can be positive when no $s_i^2$ exceeds $S_2/2$: for
\[
s=\bigl(-\tfrac34,\tfrac12,\tfrac25,\tfrac25,-\tfrac{19}{20},\tfrac25\bigr),\qquad d=\bigl(0,0,0,0,\tfrac1{20},0\bigr),
\]
the largest $s_i^2/S_2$ is $\frac{361}{878}$, $Q-H=\frac1{20000}$ and $|L_6|=\frac{167}{2500}$.

The analogue of the constant $\frac83$ of \cref{thm:four} cannot exceed~$2$: for $s=(\lambda,-\frac\lambda5,\dots,-\frac\lambda5)$ and $d=(\mu,1-\frac\lambda5,\dots,1-\frac\lambda5)$ with $0\le\mu\le1-\lambda$,
\[
Q-H=\frac{\lambda^2}{625}\bigl(500(1-\mu)-\lambda(1000-300\mu)+\lambda^2(365-35\mu)-37\lambda^3\bigr),
\]
\[
L_6=\frac{8\lambda^2(5-\lambda)^2}{625}\bigl(5(1-\mu)-\lambda(5-2\mu)+2\lambda^2\bigr),
\]
so $|L_6|/(Q-H)\to2$ as $\lambda\to0$ for each fixed $\mu<1$.

\section*{Acknowledgments and funding}
The research presented in this paper was supported by the European Research Council (ERC) under the European Union's Horizon Europe research and innovation programme (grant agreement No.~101041711), by the Simons Foundation as part of the Collaboration on the Mathematical and Scientific Foundations of Deep Learning, by Heights Labs, by Convex Nexus Capital, and by the Israel Science Foundation (grants 2258/19 and 4101/25).

This work was done by multiple AI agents, Claude (Anthropic) and Codex (OpenAI), connected to Math Brain~\cite{MathBrainRepo} and working together inside the Brain framework~\cite{BrainFramework}, a shared, persistent memory in which the agents record claims, proofs, failed attempts and methods, retrieve one another's results, coordinate who works on what, and audit each other's work. The method is described in~\cite{MathBrainPaper}.

\section*{Declaration of competing interest}
The author holds equity in Heights Labs and in Convex Nexus Capital and has a management role at Convex Nexus Capital.

\section*{Code and data availability}
The \LaTeX\ source and the programs are available at \url{https://github.com/yspennstate/poisson-binomial-centered-separation}. They verify the identities~\eqref{eq:avg4},~\eqref{eq:diff4},~\eqref{eq:Q},~\eqref{eq:affine},~\eqref{eq:regroup},~\eqref{eq:gap} and the expansion in the proof of \cref{lem:three} in exact rational arithmetic. They evaluate~\eqref{eq:average},~\eqref{eq:jumps},~\eqref{eq:exact} and~\eqref{eq:variance} by exact convolution for $1{,}335$ rational parameter vectors with $n\le20$, and~\eqref{eq:equal} for the $1{,}280$ of them with $n\ge4$. The programs \texttt{six\_summands\_identities.py} and \texttt{six\_summands\_gap.py} verify the two identities of \cref{sec:six}, the value of $b_6(g)$ and the identity in the proof of \cref{prop:six}. The program \texttt{six\_summands\_search.py} verifies the formulas of the three examples of \cref{sec:six}. It evaluates $Q-H$ and $L_6$ in floating point at $25{,}000$ feasible centered pairs with $n=6$, drawn in five classes described in the program; $17{,}132$ of them have no $s_i^2$ above $S_2/2$, $Q-H>0$ at $815$ of them ($83$ with no $s_i^2$ above $S_2/2$), and $Q-H>|L_6|$ at none of them.

\section*{Declaration of generative AI and AI-assisted technologies}
Codex (OpenAI) found the construction of \cref{sec:equal} and the argument of \cref{sec:four} and wrote the verification programs; a second Codex run checked each claim with its own code. Claude (Anthropic) found the argument of \cref{sec:six} and wrote its programs, checked the proofs and wrote the text. The question is Kontorovich's~\cite{Kontorovich}, the inequality $d_K(X,Y)\ge b_n(g)$ and its equality cases are Long's~\cite{Long}, and the proof of \cref{thm:equal} rests on Hoeffding's interval inequality~\cite{Hoeffding}. The author directed the work and is responsible for the content of the paper and for decisions concerning its dissemination.

\appendix
\section{Notation}\label{app:notation}
\begin{center}
\small
\begin{tabular}{p{0.24\textwidth}p{0.46\textwidth}p{0.2\textwidth}}
\toprule
symbol & meaning & defined in\\
\midrule
$p_i,q_i$ & the Bernoulli parameters, $p_i\ge q_i$ & \cref{sec:intro}\\
$X_i,Y_i,X,Y$ & Bernoulli$(p_i)$ and Bernoulli$(q_i)$ variables and their sums & \cref{sec:intro}\\
$n$, $m$ & number of summands; $n=2m$ when $n$ is even & \cref{sec:intro}\\
$d_i,s_i$ & $p_i-q_i$ and $p_i+q_i-1$ & \cref{sec:intro}\\
$\Delta,g,S_2$ & $\sum_id_i$, mean gap $\Delta/n$, $\sum_is_i^2$ & \cref{sec:intro}\\
$s,d$ & the vectors $(s_1,\dots,s_n)$ and $(d_1,\dots,d_n)$ & \cref{sec:intro}\\
$T$ & a quantity with $T\le d_K(X,Y)$ whose excess is considered & \cref{sec:intro}\\
$\Pp,\E$ & probability and expectation & \cref{sec:intro}\\
$S$, $c$ & a sum of independent Bernoulli variables and an integer in Hoeffding's tail comparison & \cref{sec:intro}\\
$D_k$, $d_K$ & $\Pp(X\ge k)-\Pp(Y\ge k)$; Kolmogorov distance & \cref{sec:intro}\\
$\Bin(n,p)$, $b_n(g)$ & binomial variable; Kolmogorov distance of the symmetric binomial pair & \cref{sec:intro}\\
$Q,L$ & $Q=\sum_i(d_i-\Delta/2)s_i^2$, \eqref{eq:Qdef}; $L$ the polynomial of \cref{thm:four} & \cref{sec:intro}\\
$\ell_1,\ell_2$, $\bar r$ & integer endpoints in Hoeffding's interval inequality; mean of the parameters of $S$ in \eqref{eq:hoeffding} & \cref{sec:equal}\\
$w$ & a value of $W$ & \cref{sec:intro}\\
$a$ & $1-g$ & proof of \cref{thm:equal}\\
$r_i$ & $\frac12+\frac{s_i}{2(1-g)}$ & \cref{sec:intro}\\
$R_i,Z_i,W$ & the Bernoulli$(r_i)$ and Bernoulli$(g)$ variables from which $X$ and $Y$ are built, $\sum_iR_i$ & \cref{sec:intro}\\
$\widetilde X_i,\widetilde Y_i$ & copies of $X_i,Y_i$ built from $R_i,Z_i$ & proof of \cref{thm:equal}\\
$\mathbf 1$ & indicator & proof of \cref{thm:equal}\\
$h$ & sign of $k-m$ & proof of \cref{thm:equal}\\
$\overset d=$ & equality in law & proof of \cref{thm:equal}\\
$\nu$ & an integer in $\{0,\dots,m\}$ & proof of \cref{thm:equal}\\
$\phi$ & the conditional value of $(D_m+D_{m+1})/2$ given $W=w$ & proof of \cref{thm:equal}\\
$B,B',\xi$ & $\Bin(m+j,g)$, $B+\xi$, a Bernoulli$(g)$ variable & proof of \cref{thm:equal}\\
$A_j$ & jumps of $\phi$ & proof of \cref{thm:equal}\\
$B_0$, $\delta_j$ & $\Bin(2m,\frac12)$; $\Pp(|W-m|\le j)-\Pp(|B_0-m|\le j)$ & proof of \cref{thm:equal}\\
$c^*_{2m}(g)$ & coefficient of $S_2$ obtained from~\eqref{eq:variance} & \eqref{eq:cstar}\\
$\gamma$ & a coefficient of $S_2$ in a lower bound for $(D_m+D_{m+1})/2-b_{2m}(g)$ & \cref{sec:intro}\\
$y_i$, $v_i,U,F$ & variables of \cref{lem:three}, later coordinates of $s$; quantities of \cref{lem:three} & \cref{lem:three}\\
$z$, $M$, $e_k$ & a sign vector; $1$, $0$ or $-1$ by the number of positive coordinates of $z$; elementary symmetric polynomials & proof of \cref{thm:four}\\
$x,t,u_i,\sigma$ & coordinates of $s$ and $d$; $\sum_iu_i$ & proof of \cref{thm:four}\\
$\rho,\kappa,\theta_i$ & $\sum_iy_i^2$, $(1-\rho)/2$, $(1+\rho)/2-y_i^2$ & proof of \cref{thm:four}\\
$C,G,V$ & slope and intercept of $L/x$; an auxiliary sum & proof of \cref{thm:four}\\
$H,L_6,O_j$ & the polynomials in $s$ and $d$ of the case $n=6$ & \cref{sec:six}\\
$s^{(I)}$ & $s$ with the coordinates indexed by $I$ removed & \cref{sec:six}\\
$M_6$ & the analogue of $M$ on $\{-1,1\}^6$ & \cref{sec:six}\\
\bottomrule
\end{tabular}
\end{center}
\begin{center}
\small
\begin{tabular}{p{0.24\textwidth}p{0.46\textwidth}p{0.2\textwidth}}
\toprule
term & meaning & defined in\\
\midrule
gaps, mean gap & the $d_i$; $g$ & \cref{sec:intro}\\
tail difference & $D_k$ & \cref{sec:intro}\\
feasible & $|s_i|+d_i\le1$ for every $i$ & \cref{sec:intro}\\
symmetric binomial pair & $\Bin(n,\frac{1+g}2)$, $\Bin(n,\frac{1-g}2)$ & \cref{sec:intro}\\
Poisson binomial comparison & the inequality~\eqref{eq:conj} & \cref{sec:intro}\\
excess & $T-b_n(g)$ & \cref{sec:intro}\\
centered & $\sum_i(p_i+q_i)=n$, that is $\sum_is_i=0$ & \cref{sec:intro}\\
$\eta;\beta,\tau;\lambda,\mu$ & a point of $[0,1]^4$; parameters of two families of pairs & \cref{sec:six}\\
\bottomrule
\end{tabular}
\end{center}

\begin{thebibliography}{99}

\bibitem{Gleser} L.~J.~Gleser, On the distribution of the number of successes in independent trials, \emph{Ann. Probab.} \textbf{3} (1975), no.~1, 182--188, doi:10.1214/aop/1176996461.

\bibitem{Hoeffding} W.~Hoeffding, On the distribution of the number of successes in independent trials, \emph{Ann. Math. Statist.} \textbf{27} (1956), no.~3, 713--721, doi:10.1214/aoms/1177728178.

\bibitem{Kontorovich} A.~Kontorovich, Poisson binomial conjecture, MathOverflow question 477402 (2024), \url{https://mathoverflow.net/q/477402}.

\bibitem{KontorovichHom} A.~Kontorovich, A homogenization principle for total variation, arXiv:2604.03882 (2026).

\bibitem{KontorovichTV} A.~Kontorovich, TV homogenization inequalities, arXiv:2601.04079 (2026).

\bibitem{KontorovichTens} A.~Kontorovich, On the tensorization of the variational distance, Electron. Commun. Probab. \textbf{30} (2025), paper no.~32, doi:10.1214/25-ECP680.

\bibitem{Long} C.~D.~Long, The exact Poisson-binomial comparison in every dimension, manuscript dated 13 September 2026, with a Lean formalization, \url{https://github.com/long-mathematics/poisson-binomial-comparison}, commit \texttt{b34cde6}.

\bibitem{Sawin} W.~Sawin, answer to MathOverflow question 477402 (2024), \url{https://mathoverflow.net/a/479497}.

\bibitem{MathBrainRepo} Y.~Shmalo, \emph{Math Brain: public snapshot, code and tests}, 2026, \url{https://github.com/yspennstate/math-brain-public}.

\bibitem{MathBrainPaper} Y.~Shmalo, \emph{Math Brain: shared, auditable memory for multi-agent mathematical research}, October 2026, \url{https://github.com/yspennstate/math-brain-public/tree/main/papers/math-brain-2026-10-24}.

\bibitem{BrainFramework} Y.~Shmalo, \emph{Brain Framework}, public software repository, 2026, \url{https://github.com/yspennstate/brain-framework-public}.

\bibitem{Steinke} T.~Steinke, answer to MathOverflow question 477402 (2024), \url{https://mathoverflow.net/a/477441}.

\bibitem{TangTang} W.~Tang and F.~Tang, The Poisson binomial distribution---old \& new, \emph{Statist. Sci.} \textbf{38} (2023), no.~1, 108--119, doi:10.1214/22-STS852.

\end{thebibliography}
\end{document}
